\chapter{FLT - Fried Lettuce and Tomato Sandwich?}

In two episodes (which first aired on October 30, 1995
and September 20, 1998), we see a statement of the form 
$a^{12}+b^{12}=c^{12}$, where $a$, $b$, and $c$ are natural numbers.
These equations are the following.
\begin{equation}
1782^{12}+1841^{12}=1922^{12}\label{funnyeqn1}
\end{equation}

\begin{equation}
3987^{12}+4365^{12}=4472^{12}\label{funnyeqn2}
\end{equation}

This may call to mind a little-known theorem.

\begin{thm}[Fermat's Last Theorem, or FLT]\label{flt}
The equation $x^n+y^n=z^n$ has no non-zero integer solutions for $n>2$.
\end{thm}

\begin{proof}
I have a marvellous proof of this, but the grant funding to produce this
manual is too small for me to present it here. However, interested 
readers may wish to see \cite{AW1995}.\\
\end{proof}

A consequence of this theorem, of course, is that equations (\ref{funnyeqn1})
and (\ref{funnyeqn2}) must be incorrect. But how might one determine
this without recourse to FLT? Well, if one had lots of spare time late
at night (as you very well may, dear reader), one could just go ahead and
multiply them out. This author, on the other hand, has a very active
social life. So, we instead recall the following.

\begin{lem}\label{integerproducts}
The product of any two odd integers is an odd integer. The product of any 
two even integers is an even integer.
\end{lem}

\begin{proof}
Look in your notes or text from Math Camp. This is one of many similar
results we proved back in those good old days.\\
\end{proof}

\begin{cor}\label{oddtopower}
An odd integer raised to any positive integer power is an odd integer.
That is, for odd $z\in\mathbb{Z}$, $z^m$ is odd for any $m\in\mathbb{N}$.
\end{cor}

\begin{proof}
By induction. Suppose $z$ is any odd integer.

\emph{Base case}: Since $z$ is odd, we know that $z^1=z$ is odd.

\emph{Inductive step}: Suppose $z^n$ is odd for some $n\in\mathbb{N}$ 
(our inductive hypothesis). We must show that $z^{n+1}$ is also odd.

Well,... $z^{n+1}$=$z^n\cdot z$ by basic properties of exponents.

By our inductive hypothesis we know that $z^n$ is odd, and we began
by assuming that $z$ is odd. So, $z^n\cdot z$ is the product of two odd 
integers and is therefore odd by Lemma (\ref{integerproducts}).\\
\end{proof}

\begin{cor}\label{eventopower}
An even integer raised to any positive integer power is an even integer.
\end{cor}

\begin{proof}
This proof is nearly identical to the one above. To reduce the printing
cost of this thesis and conserve some of the planet's natural resources,
we omit it.\\
\end{proof}

We now look back to equation (\ref{funnyeqn1}), $1782^{12}+1841^{12}=1922^{12}$,
and see that we have an two even integers, 1782 and 1922, raised to 
positive integer powers; thus, by Corollary \ref{eventopower}, each of 
the results is even. Similarly, by Corollary \ref{oddtopower}, $1841^{12}$
is odd.\\

Now all that remains is to invoke one of the most powerful theorems known
to mankind.

\begin{thm}\label{awesometheorem}
The sum of an even integer and an odd integer is an odd integer.
\end{thm}

\begin{proof}
The proof of this theorem is so far beyond the scope of this thesis as
to be nearly incomprehensible to a person of your intellect. Some of the 
greatest math students in the history of Colby College are currently 
devoting all their talents and energies to a proof, even to the extent 
that they have had to reduce their daily hours spent playing Dungeons 
and Dragons from sixteen to just two. Fortunately, they are unemployed
and living in their mothers' basements, so they have little else to do.\\
\end{proof}

By applying Theorem \ref{awesometheorem} to the left-hand side of
Equation (\ref{funnyeqn1}), we see that this left-hand side must
be odd. But we already found above that the right-hand side is even.
Thus, this equation cannot be true.\\

Alas, this same argument cannot be applied to Equation (\ref{funnyeqn2}).
[Why not?] However, by using a different base for our modular arithmetic,
we can arrive at the same conclusion. This, my loyal and persevering
reader, is left on your shoulders. So put away your video games and
your crazy rock and roll records and get to work, you slacker.
